% Main-text candidate. Environments proposition/theorem/corollary/remark required.
% All identifiers start sc: to avoid collisions with matrix-agent contributions.

\subsection{From one deterministic edge to several stochastic diagnostics}
For deterministic $L(x)=hx^2/2$ with $h>0$, the entire orbit is
\begin{equation}
 x_t=(1-\eta h)^t x_0.
 \label{eq:deterministic}
\end{equation}
The condition $0<\eta h<2$ contracts every nonzero start to zero. At $\eta h=2$ the sign alternates without decay; above two the magnitude diverges. A fixed quadratic therefore cannot explain bounded large-step fluctuations after its expanding boundary. Random curvature changes the product law, and nonlinear curvature can change the orbit's return.

On the scalar sampled quadratic
\(\ell_B(x)=h_Bx^2/2-\xi_Bx\), SGD is exactly
\begin{equation}
 X_{t+1}=A_{t+1}X_t+\eta\xi_{t+1},\qquad A_B=1-\eta h_B.
 \label{sc:quadratic}
\end{equation}
Two trajectories using the same sampled batches differ by
\(\Delta_{t+1}=A_{t+1}\Delta_t\).
Write \(M=|A_B|\), so moments always use the absolute multiplier, and define
\begin{equation}
 m(p)=\mathbb E M^p,\quad
 \gamma=\mathbb E\log M,\quad
 H=\frac1{m(-1)}.
 \label{sc:diagnostics}
\end{equation}
Here \(H\) is the harmonic gain; its inverse \(m(-1)\) is the negative-moment
coefficient. We only assign a finite inverse-moment estimate this meaning
when the population inverse moment exists. For constant curvature and $p\ne0$,
\(m(p)^{1/p}=M\), \(\gamma=\log M\), and \(H=M\): the boundaries all
coincide at \(|1-\eta h|=1\). Curvature randomness separates them.

\begin{proposition}[Scalar moment ordering]\label{sc:ordering}
Let \(M>0\), \(\mathbb E|\log M|<\infty\), and suppose the displayed moments
are finite. Then
\begin{equation}
 \log H\ \le\ \gamma\ \le\ \log m(1)\ \le\ \tfrac12\log m(2).
 \label{sc:ladder}
\end{equation}
For iid multipliers and \(\Delta_0\ne0\),
\(t^{-1}\log|\Delta_t/\Delta_0|\to\gamma\) almost surely and
\(\mathbb E|\Delta_t/\Delta_0|^p=m(p)^t\).
Consequently \(H>1\) implies \(\gamma>0\). The converse need not hold.
\end{proposition}

The use of absolute values matters: $m(1)$ is a mean-absolute multiplier, not the signed mean $\mathbb E a_B$. For $p>0$, $m(p)<1$ means contraction of a $p$th perturbation moment. For $p<0$ it means contraction of an \emph{inverse distance}, so it measures repulsion from zero, rather than contraction toward zero. The logarithmic criterion lies between these orders. We use the term harmonic expansion for $H>1$ while keeping this direction of the inequality explicit.

A concrete two-curvature law already separates the references. Let $a_B=-1/2$ with probability $u$ and $a_B=-2$ otherwise; both are realizable with positive quadratic curvatures. Then
\begin{equation}
 \gamma=(1-2u)\log2,\qquad H=(1/2+3u/2)^{-1}.
 \label{eq:two-curvatures}
\end{equation}
For $1/3<u<1/2$, the product is logarithmically expanding but $H<1$. For $u<1/3$, both diagnostics expand. No dimension or gradient-alignment assumption is needed for this separation.

The classical Kesten--Goldie result concerns the \emph{affine} recursion
\eqref{sc:quadratic}: with \(\gamma<0\), nondegenerate additive forcing, and
the usual renewal assumptions, its stationary upper tail has exponent
\(\kappa>0\) solving \(m(\kappa)=1\)
\citep{kesten1973random,goldie1991implicit}.
This is a baseline, not a theorem of post-expansion confinement. In
particular, the nondegenerate Gaussian affine model used in
Appendix~\ref{sc:proofs} has no invariant probability when \(\gamma>0\).
An additional return mechanism is needed.

\subsection{Nonlinear return changes an upper-tail question into a central-mass question}
An \emph{invariant law} $\pi$ satisfies $X_0\sim\pi\Rightarrow X_t\sim\pi$ at every time. It is \emph{attracting} when other initial laws converge to it in distribution. These statements concern the law of a state; they need not imply contraction of paths using the same random updates.

The simplest exact return model is saturated amplitude:
\begin{equation}
 R_{t+1}=\min\{r_*,M_{t+1}R_t\},\qquad 0<R_t\le r_*.
 \label{sc:saturation}
\end{equation}
Here $r_*>0$ is the outer return radius. The model separates local multiplicative motion from return at this scale.
It is an explicit model, not an identity for a globally quadratic loss.

\begin{theorem}[An inverse-moment law after local expansion]\label{sc:central-law}
Suppose \(\mathbb E|\log M|<\infty\), \(\gamma>0\),
\(\log M\) is nonarithmetic, and
\(m(-\alpha)=1\) for some \(\alpha>0\). Assume \(m(-q)\) is finite in a
neighborhood of \(\alpha\) and
\(0<-\mathbb E[M^{-\alpha}\log M]<\infty\).
Then \eqref{sc:saturation} has a unique invariant probability on
\((0,r_*]\), and for some \(C\in(0,\infty)\),
\begin{equation}
 \mathbb P_\pi(R\le r)\sim C(r/r_*)^\alpha,\qquad r\downarrow0.
 \label{sc:smallball}
\end{equation}
If, additionally, \(\log M\) has a bounded continuous density of compact
support, then
\(p_R(r)\sim C\alpha r_*^{-\alpha}r^{\alpha-1}\).
When the negative-moment curve is strictly convex and finite through
\(1\) and \(\alpha\),
\begin{equation}
 H<1\Longleftrightarrow\alpha<1,\qquad
 H=1\Longleftrightarrow\alpha=1,\qquad
 H>1\Longleftrightarrow\alpha>1.
 \label{sc:shape-ladder}
\end{equation}
\end{theorem}

Thus positive local log growth can coexist with a central density cusp
\((0<\alpha<1)\); the harmonic crossing separates this cusp from central
depletion \((\alpha>1)\) in this return model. For a reflection-symmetric
signed coordinate, the same power occurs on each side. Central depletion
is a local shape conclusion. Exactly two modes additionally requires
control of the outer density, and arbitrary reinjection can reverse the
conclusion (Appendix~\ref{sc:counterexample}).

\paragraph{Proof idea.}
The change of variable \(W=\log(r_*/R)\) gives
\(W_{t+1}=\max\{0,W_t-\log M_{t+1}\}\).
This is a reflected random walk with negative mean increment.
Its invariant law is the supremum of the unreflected walk. The nonarithmetic assumption rules out log-radius lattice effects. Tilting the
increment law by \(M^{-\alpha}\) reverses its drift; the overshoot renewal
theorem gives an exponential tail in \(W\), hence the power
\eqref{sc:smallball} in \(R\). The root therefore has the opposite sign
from the Kesten upper-tail root.
